\section{Reference Shader Implementations}
\label{app:shaders}

The body gives the kernels whose \emph{form} carries an argument: the stable
characteristic curve (Listing~\ref{lst:charcurve}), the interlayer combine, the
grain field generator (Listing~\ref{lst:grain}), the fused pointwise chain, and
tetrahedral interpolation (Listing~\ref{lst:tetra}). This appendix supplies the
remainder needed to reproduce the pipeline, in the order they execute.

All listings are Metal Shading Language 3.1 unless noted, compile with
\texttt{-ffast-math} disabled (Section~\ref{sec:precision}), and assume the
threadgroup sizing and bounds convention established in
Section~\ref{sec:metal}: every kernel begins with a grid guard, because
dispatch sizes are rounded up to threadgroup multiples and a missing guard
corrupts the row below the image in a way that is invisible until export.

\subsection{Host--Shader Parameter ABI}

One header is shared between Swift and Metal, so a parameter added on one side
fails to compile on the other rather than silently misaligning. Alignment is
explicit: Metal's \texttt{float3} is 16-byte aligned, and a Swift
\texttt{SIMD3<Float>} matches it, but a Swift tuple does not.

\begin{lstlisting}[language=Metal, caption={Shared parameter block (\texttt{EmulsionABI.h}).}, label={lst:abi}]
#ifndef EMULSION_ABI_H
#define EMULSION_ABI_H
#include <simd/simd.h>

// Per-record characteristic curve parameters (softplus density model).
typedef struct {
    simd_float3 dMin;      // base + fog, per record
    simd_float3 deltaD;    // density range
    simd_float3 gamma;     // straight-line gamma (negative for reversal)
    simd_float3 x0;        // extrapolated speed point, log10 lx-s
    simd_float3 kappaT;    // toe softness
    simd_float3 kappaS;    // shoulder softness
    simd_float3 maskDepl;  // m, mask depletion
} EMCurveParams;

typedef struct {
    simd_float3x3 dirMatrix;   // A, DIR coupling matrix
    simd_float2   dirSigmaPx;  // (sigma1, sigma2) in render pixels
    float         dirW2;       // long-scale weight
    float         dirMu;       // activity exponent
} EMInterlayerParams;

typedef struct {
    simd_float3   sigmaPx;     // grain kernel sigma per record, pixels
    simd_float3x3 cholesky;    // L, chroma correlation factor
    simd_float2   nu;          // granularity shape exponents
    float         amount;      // user grain amount g
    float         chi;         // clustering fraction
    float         lambdaPx;    // expected grains per pixel
    uint          seed;        // frame seed
} EMGrainParams;

typedef struct {
    simd_float3 aimLogL;       // log10 L_aim, computed aim balance
    simd_int3   points;        // printer lights p_c, integer
    int         masterPoints;  // print density
    simd_float3x3 C;           // printing density matrix
    EMCurveParams print;       // print stock curve
    float       silverRho;     // silver retention fraction
    float       surroundGamma; // viewing condition exponent
} EMPrintParams;

typedef struct {
    EMCurveParams      negative;
    EMInterlayerParams interlayer;
    EMGrainParams      grain;
    EMPrintParams      print;
    simd_float3x3      mIn;        // P3 -> ACEScg input transform
    simd_float3x3      mWhite;     // collapsed von Kries adaptation
    float              anchorGain; // exposure anchoring gain
    uint               qualityTier;
} EMFrameParams;
#endif
\end{lstlisting}

\texttt{EMFrameParams} is a single 3-buffer upload of under 700 bytes, written
once per frame into a triple-buffered ring. Nothing in the pipeline reads a
parameter from any other source, which is what makes the render a pure function
of \texttt{(source texture, EMFrameParams)} and therefore makes the golden-image
tests of Section~\ref{sec:validation} meaningful.

\subsection{Input Transform, Anchoring, and White Balance}

The three host-computed matrices collapse into one, so the entire color
management front end costs a single multiply. This is worth stating explicitly
because it is easy to implement as three passes and never notice the cost.

\begin{lstlisting}[language=Metal, caption={Fused input transform.}]
kernel void em_input_transform(
    texture2d<float, access::read>  src  [[texture(0)]],
    texture2d<float, access::write> dst  [[texture(1)]],
    constant EMFrameParams&         P    [[buffer(0)]],
    uint2 gid [[thread_position_in_grid]])
{
    if (gid.x >= dst.get_width() || gid.y >= dst.get_height()) return;

    float3 e = src.read(gid).rgb;          // extended linear Display P3

    // Host has premultiplied M_in * M_white; one matrix, not two.
    e = P.mWhite * (P.mIn * e);

    // Exposure anchoring: anchorGain folds 2^eps * 0.18 / g_cal.
    e *= P.anchorGain;

    // The pipeline takes log10 next; a non-positive channel is not an error
    // to clamp away silently, it is a signal that the working space is too
    // small. AP1 makes this rare; we floor at a value below any real signal
    // and count occurrences in the debug build.
    e = max(e, float3(1e-7));

    dst.write(float4(e, 1.0), gid);
}
\end{lstlisting}

\subsection{Separable Recursive Gaussian}
\label{app:recursive}

The interlayer diffusion blurs and the grain kernel convolution use a separable
recursive filter \cite{deriche1993,young1995}, whose cost is independent of
$\sigma$. This is what makes the interlayer long scale --- $6\,\mu$m at the film
plane, tens of pixels at export resolution --- affordable at all. The
implementation is a third-order forward/backward pass per axis; coefficients are
computed on the host from $\sigma$ and uploaded.

Halation does \emph{not} use this filter. As established in
Section~\ref{sec:recursive-blur}, a separable IIR realization of the exponential
PSF is anisotropic in a way that is plainly visible on point highlights, so
halation is evaluated instead as the fitted Gaussian mixture over mip levels,
each term of which is small enough for direct separable convolution. The
distinction matters: the recursive filter is correct where the target kernel is
Gaussian and the radius is large, and wrong where the target kernel is an
isotropic exponential.

\begin{lstlisting}[language=Metal, caption={Recursive Gaussian, one axis, one row per thread.}, label={lst:recursive}]
struct EMRecursiveCoeffs { float b0, b1, b2, b3, B; };

kernel void em_recursive_blur_h(
    texture2d<float, access::read>  src [[texture(0)]],
    texture2d<float, access::write> dst [[texture(1)]],
    constant EMRecursiveCoeffs&     k   [[buffer(0)]],
    uint gy [[thread_position_in_grid]])
{
    const uint W = src.get_width();
    if (gy >= src.get_height()) return;

    // Forward pass. Boundary is initialized to the edge value rather than
    // zero: a zero boundary produces a dark rim that survives into the
    // print and reads as a vignette the user did not ask for.
    float3 p1, p2, p3;
    p1 = p2 = p3 = src.read(uint2(0, gy)).rgb;

    for (uint x = 0; x < W; ++x) {
        float3 s = src.read(uint2(x, gy)).rgb;
        float3 o = k.B * s + k.b1 * p1 + k.b2 * p2 + k.b3 * p3;
        p3 = p2; p2 = p1; p1 = o;
        dst.write(float4(o, 1.0), uint2(x, gy));
    }

    // Backward pass, in place.
    p1 = p2 = p3 = dst.read(uint2(W - 1, gy)).rgb;
    for (int x = int(W) - 1; x >= 0; --x) {
        float3 s = dst.read(uint2(uint(x), gy)).rgb;
        float3 o = k.B * s + k.b1 * p1 + k.b2 * p2 + k.b3 * p3;
        p3 = p2; p2 = p1; p1 = o;
        dst.write(float4(o, 1.0), uint2(uint(x), gy));
    }
}
\end{lstlisting}

The vertical pass is the transpose of this kernel. Because a row-per-thread
dispatch has poor occupancy at small image heights, the graph substitutes the
threadgroup-staged two-pass separable convolution of Section~\ref{sec:metal}
below $\sigma = 2.5$ pixels, where the direct filter is cheaper anyway; the
crossover was measured, not assumed (Section~\ref{sec:performance}).

\subsection{Halation Composite}

The four-term Gaussian mixture of Table~\ref{tab:gaussmix} is evaluated by
blurring a thresholded source at four mip levels and accumulating. The
threshold is soft, because a hard threshold produces a visible contour in a
sky at the point where the source crosses it.

\begin{lstlisting}[language=Metal, caption={Halation source extraction and composite.}]
kernel void em_halation_source(
    texture2d<float, access::read>  lin [[texture(0)]],
    texture2d<float, access::write> out [[texture(1)]],
    constant EMHalationParams&      H   [[buffer(0)]],
    uint2 gid [[thread_position_in_grid]])
{
    if (gid.x >= out.get_width() || gid.y >= out.get_height()) return;
    float3 e = lin.read(gid).rgb;

    // Soft threshold: smoothstep over a knee of width kappa_h around tau.
    float3 t = smoothstep(float3(H.tau - H.kappaH),
                          float3(H.tau + H.kappaH), e);
    float3 s = pow(max(e - H.tau, float3(0.0)), H.q) * t;

    // Channel bias beta applies at the source, not at the composite, so
    // that a red-biased halo also scatters over the red channel's longer
    // length -- the two are the same physical fact.
    out.write(float4(s * H.beta, 1.0), gid);
}

kernel void em_halation_composite(
    texture2d<float, access::read>       base [[texture(0)]],
    array<texture2d<float, access::sample>, 4> mips [[texture(1)]],
    texture2d<float, access::write>      dst  [[texture(5)]],
    constant EMHalationParams&           H    [[buffer(0)]],
    sampler                              samp [[sampler(0)]],
    uint2 gid [[thread_position_in_grid]])
{
    if (gid.x >= dst.get_width() || gid.y >= dst.get_height()) return;
    float2 uv = (float2(gid) + 0.5) / float2(dst.get_width(), dst.get_height());

    float3 halo = float3(0.0);
    for (uint j = 0; j < 4; ++j) {
        halo += H.w[j] * mips[j].sample(samp, uv).rgb;
    }
    // Base-reflection ring: a second, wider term weighted by omega, offset
    // to the minimum radius set by base thickness. Near zero for modern
    // stocks; the .xr profiles are where it earns its place.
    halo += H.omega * mips[3].sample(samp, uv).rgb;

    float3 e = base.read(gid).rgb + H.alpha * halo;
    dst.write(float4(e, 1.0), gid);
}
\end{lstlisting}

Halation is added in \emph{linear exposure}, before the characteristic curve,
because that is where the scattered light physically lands. Adding it to
density, or to the print, produces a glow that does not compress in the
negative's shoulder --- which is the tell that separates a simulated halo from
a bloom filter.

\subsection{Poisson-Mode Grain}

Listing~\ref{lst:grain} gives the Gaussian mode used at full resolution.
Below $\Lambda = 10$ expected grains per pixel the central limit theorem fails
and individual grains must be realized.

\begin{lstlisting}[language=Metal, caption={Poisson-mode grain realization.}]
inline uint em_hash(uint3 v) {        // shared with the Gaussian-mode kernel
    uint h = v.x * 0x9E3779B9u ^ v.y * 0x85EBCA6Bu ^ v.z * 0xC2B2AE35u;
    h ^= h >> 15; h *= 0x2545F491u; h ^= h >> 13;
    return h;
}
inline float em_unit(uint h) { return float(h & 0x00FFFFFFu) / 16777216.0f; }

kernel void em_grain_poisson(
    texture2d<float, access::read>  dens [[texture(0)]],
    texture2d<float, access::write> fld  [[texture(1)]],
    constant EMGrainParams&         G    [[buffer(0)]],
    uint2 gid [[thread_position_in_grid]])
{
    if (gid.x >= fld.get_width() || gid.y >= fld.get_height()) return;

    // Lattice cell in FILM coordinates, not pixel coordinates: this is what
    // makes preview and export the same field at two sampling rates.
    uint3 cell = uint3(gid.x, gid.y, G.seed);

    float3 acc = float3(0.0);
    // 3x3 neighbourhood covers the kernel support at Poisson-mode sigmas.
    for (int dy = -1; dy <= 1; ++dy) {
        for (int dx = -1; dx <= 1; ++dx) {
            uint  h  = em_hash(cell + uint3(uint(dx + 1), uint(dy + 1), 0u));
            // Knuth's method; valid because lambda < 10 by construction.
            float L  = exp(-G.lambdaPx);
            float p  = 1.0; uint n = 0u;
            do { p *= em_unit(em_hash(uint3(h, n, 0u))); ++n; } while (p > L);
            n -= 1u;

            for (uint i = 0; i < n; ++i) {
                uint  hi = em_hash(uint3(h, i, 1u));
                float ox = em_unit(hi)            + float(dx);
                float oy = em_unit(hi * 747796405u) + float(dy);
                float r2 = ox * ox + oy * oy;
                acc += exp(-r2 / (2.0 * G.sigmaPx * G.sigmaPx));
            }
        }
    }
    // Normalize to zero mean, unit variance; the density-dependent
    // scaling sigma_D(D) is applied in the combine pass.
    float3 mean = G.lambdaPx * 2.0 * M_PI_F * G.sigmaPx * G.sigmaPx;
    fld.write(float4((acc - mean) * rsqrt(mean), 1.0), gid);
}
\end{lstlisting}

\subsection{Print Transfer, Silver Retention, and Display}

The print stage is a single fused kernel: crosstalk, exposure, curve,
retention, and display normalization touch each pixel once. It is the last
place in the pipeline where a texture round trip would be pure waste, since
every operation is pointwise given the negative density.

\begin{lstlisting}[language=Metal, caption={Fused print transfer and display conversion.}, label={lst:print}]
inline float3 em_softplus(float3 u, float3 a) {   // numerically stable form
    return max(u, 0.0) + a * log1p(exp(-abs(u) / a));
}

kernel void em_print_transfer(
    texture2d<float, access::read>  negD [[texture(0)]],
    texture2d<float, access::write> dst  [[texture(1)]],
    constant EMFrameParams&         P    [[buffer(0)]],
    uint2 gid [[thread_position_in_grid]])
{
    if (gid.x >= dst.get_width() || gid.y >= dst.get_height()) return;
    float3 D = negD.read(gid).rgb;

    // Printing density: unwanted absorptions of the negative's dyes.
    float3 Deff = P.print.C * D;

    // Printer lights as a 0.025-per-point offset from the computed aim.
    float3 pts = float3(P.print.points) * 0.025
               + float(P.print.masterPoints) * 0.025;
    float3 logE = P.print.aimLogL + pts - Deff;

    // Print characteristic curve: identical form to the negative.
    const EMCurveParams k = P.print.print;
    float3 u  = k.gamma * (logE - k.x0);
    float3 Dp = k.dMin + em_softplus(u, k.kappaT)
                       - em_softplus(u - k.deltaD, k.kappaS);

    // Silver retention: a neutral density proportional to total
    // development, added to all three records.
    if (P.print.silverRho > 0.0) {
        float bar = dot((Dp - k.dMin) / k.deltaD, float3(1.0 / 3.0)) * 0.90;
        Dp += P.print.silverRho * bar;
    }

    // Print density to relative luminance. The Dmax subtraction
    // is not cosmetic: without it the print's finite black compounds with
    // the panel's and the image looks washed out on OLED.
    float3 dMax = k.dMin + k.deltaD;
    float3 num  = exp10(-Dp)     - exp10(-dMax);
    float3 den  = exp10(-k.dMin) - exp10(-dMax);
    float3 Y    = clamp(num / den, 0.0, 1.0);

    if (P.print.surroundGamma != 1.0) Y = pow(Y, P.print.surroundGamma);

    dst.write(float4(Y, 1.0), gid);   // linear, print primaries
}
\end{lstlisting}

\subsection{Aging Artifacts}

Dust and scratches are generated procedurally in the density domain with the
polarity of Table~\ref{tab:aging}, from the same hash as the grain field but a
different salt, so a frame's blemishes are as reproducible as its grain.

\begin{lstlisting}[language=Metal, caption={Dust and scratch overlay in negative density.}]
kernel void em_aging(
    texture2d<float, access::read_write> D [[texture(0)]],
    constant EMAgingParams&              A [[buffer(0)]],
    uint2 gid [[thread_position_in_grid]])
{
    if (gid.x >= D.get_width() || gid.y >= D.get_height()) return;
    float3 d = D.read(gid).rgb;

    // Dust: sparse opaque particles. Positive density in the negative,
    // therefore WHITE in the print -- the opposite of what an overlay
    // texture applied to the final image would give, and the reason those
    // always look wrong.
    uint  h = em_hash(uint3(gid, A.seed ^ 0x5BD1u));
    float u = em_unit(h);
    if (u < A.dustDensity) {
        float r = em_unit(h * 2654435761u);
        d += A.dustOpacity * (0.4 + 0.6 * r);
    }

    // Emulsion scratch: a vertical run of removed emulsion, negative
    // density, printing dark. Column-coherent by hashing x only.
    uint  hc = em_hash(uint3(gid.x, 0u, A.seed ^ 0x1D3Bu));
    if (em_unit(hc) < A.scratchDensity) {
        float depth = 0.3 + 0.7 * em_unit(hc * 40503u);
        d -= A.scratchStrength * depth;
    }

    D.write(float4(max(d, 0.0), 1.0), gid);
}
\end{lstlisting}

\subsection{Pointwise Chain Bake}

The LUT bake of Section~\ref{sec:lutbake} runs the entire pointwise chain ---
layer balance, characteristic curve, mask depletion, printing density, print
curve, silver retention, neutral axis, output transform --- over the $45^3$
lattice in a single dispatch of 91\,125 threads. It is the same code path as
the per-pixel chain, invoked with lattice coordinates instead of texture reads,
which is the only way to guarantee that the baked LUT and the direct evaluation
agree; a separately written bake kernel drifts from the real chain within two
releases, and the drift shows up as a preview that does not match the export.

\begin{lstlisting}[language=Metal, caption={LUT bake driver.}, label={lst:bake}]
kernel void em_bake_lut(
    texture3d<half, access::write> lut [[texture(0)]],
    constant EMFrameParams&        P   [[buffer(0)]],
    uint3 gid [[thread_position_in_grid]])
{
    const uint N = 45;
    if (gid.x >= N || gid.y >= N || gid.z >= N) return;

    // Inverse shaper: the lattice is uniform in log exposure over [-4, +2].
    // A lattice uniform in linear exposure would spend almost all of its
    // nodes above middle gray and band the shadows.
    float3 t    = float3(gid) / float(N - 1);
    float3 logE = t * 6.0 - 4.0;

    float3 Y = em_pointwise_chain(logE, P);   // shared with the direct path
    lut.write(half4(half3(Y), 1.0h), gid);
}
\end{lstlisting}

The equality of the two paths is asserted by a test rather than by inspection:
the bake is run, the direct chain is evaluated at 20\,000 random points, and the
maximum density difference must remain below $7\times10^{-3}$ --- the
tetrahedral interpolation bound derived in Section~\ref{sec:lutbake}.
